\chapter{Thị giác Đa phương thức và các Ứng dụng Tiên phong}
\label{chap:multimodal_and_advanced_applications}
%----------------------------------------------------------------------------------------
%	INTRODUCTION
%----------------------------------------------------------------------------------------

Hành trình của chúng ta qua thế giới thị giác máy tính đã đi từ những khối xây dựng cơ bản nhất—các pixel và cạnh—đến những kiến trúc nơ-ron phức tạp có khả năng học, nhận dạng, và thậm chí là sáng tạo. Chúng ta đã dạy cho máy móc cách "nhìn". Nhưng "nhìn" thôi là chưa đủ. Trí thông minh thực thụ không tồn tại trong một không gian biệt lập; nó được định hình bởi sự tương tác, bối cảnh, và khả năng kết nối các luồng thông tin khác nhau. Con người chúng ta không chỉ nhìn thấy một con chó; chúng ta có thể nói "đó là một chú chó golden retriever đang chơi đùa trong công viên" và hiểu được những hàm ý đằng sau đó.

Chương cuối cùng của phần lý thuyết này sẽ đưa chúng ta đến biên giới quan trọng nhất của thị giác máy tính hiện đại, nơi những bức tường ngăn cách giữa các giác quan kỹ thuật số đang sụp đổ: \textbf{Thị giác Đa phương thức (Multimodal Vision)}. Đây là bước nhảy vọt từ việc \textit{nhận thức} (perceiving) một cách thụ động sang \textit{lý luận} (reasoning) và \textit{tương tác} (interacting) một cách chủ động.

Trọng tâm của chương này xuất phát từ một câu hỏi đơn giản nhưng sâu sắc: Làm thế nào để dạy máy tính "nói" về những gì nó "thấy"? Làm thế nào để nó có thể hiểu được một câu lệnh bằng ngôn ngữ tự nhiên và liên kết nó với các đối tượng trong một bức ảnh hay một đoạn video? Đây là kỷ nguyên của sự hợp nhất giữa thị giác và ngôn ngữ, nơi các pixel, từ ngữ, và thậm chí cả các chuỗi hành động được đưa vào cùng một không gian ý nghĩa.

Chúng ta sẽ khám phá sự tiến hóa của các hệ thống đa phương thức này:
\begin{itemize}
	\item Chúng ta sẽ bắt đầu với cuộc cách mạng \textbf{CLIP}, mô hình đã tạo ra cầu nối đầu tiên vững chắc giữa không gian hình ảnh và không gian văn bản, mở ra kỷ nguyên của khả năng hiểu thị giác \textbf{zero-shot} và đặt nền móng cho gần như mọi mô hình đa phương thức sau này.
	
	\item Từ đó, chúng ta sẽ chứng kiến sự trỗi dậy của các \textbf{Mô hình Ngôn ngữ Lớn Đa phương thức (LVLMs)} như GPT-4o, LLaVA, và Gemini. Đây không còn là các mô hình thị giác đơn thuần; chúng là những "bộ não" kỹ thuật số có khả năng "nhìn", "đọc", và "trò chuyện" về thế giới, tạo ra một hình thức tương tác người-máy hoàn toàn mới.
	
	\item Chúng ta sẽ mở rộng khái niệm này sang chiều thời gian với các \textbf{Mô hình Ngôn ngữ-Video (Video-Language Models)}, khám phá cách các hệ thống AI có thể hiểu và mô tả các hành động, sự kiện và cốt truyện diễn ra trong các video.
	
	\item Cuối cùng, chúng ta sẽ xem xét cách những năng lực nhận thức đa phương thức này được triển khai vào thế giới thực. Chúng ta sẽ tìm hiểu về vai trò của thị giác trong các lĩnh vực tiên phong như \textbf{Robot học (Robotics)}, \textbf{Thực tế Tăng cường/Thực tế ảo (AR/VR)}, và các thách thức trong việc triển khai các mô hình này trên các \textbf{thiết bị biên (Edge Deployment)}.
\end{itemize}

Chương này là sự hội tụ của tất cả các ý tưởng mà chúng ta đã xây dựng: CNN và Transformer cung cấp "đôi mắt", học tự giám sát và các mô hình nền tảng cung cấp "bộ não" có khả năng khái quát hóa, và giờ đây, sự kết hợp với ngôn ngữ sẽ mang lại "tiếng nói" và khả năng tương tác. Chào mừng bạn đến với chương cuối cùng của cuộc hành trình lý thuyết, nơi chúng ta sẽ chứng kiến sự ra đời của một thế hệ AI không chỉ nhìn, mà còn hiểu, giao tiếp, và tương tác với thế giới của chúng ta một cách sâu sắc và đa chiều.